%--------------------------------------------------------------------%
% PETUNJUK PENGISIAN
% Halaman ini otomatis hilang bila \petunjukfalse di thesis.tex.
%--------------------------------------------------------------------%
\chapter*{PETUNJUK PENGISIAN}
\thispagestyle{capstone}
{\itshape Halaman ini dan halaman Peta Kewajiban dihapus dari naskah final (atur \texttt{\textbackslash petunjukfalse} di \texttt{thesis.tex}). Isi setiap bab dan subbab dalam bentuk paragraf. Teks miring berawalan ``Petunjuk:'' menjelaskan isi yang diharapkan dan otomatis hilang saat saklar petunjuk dimatikan; teks abu-abu dalam kurung siku diganti dengan isi Anda.

Tabel hanya dipakai untuk menyajikan data yang memang berbentuk tabel, diberi nomor dan judul (Tabel 2.1 di atas tabel), dan dirujuk dalam narasi. Jangan menulis isi bab atau subbab di dalam tabel atau kotak.\par}

\begin{itemize}
	\item Judul penelitian maksimal 16 kata dan ditulis sama persis di sampul, lembar pengesahan, pernyataan, dan presentasi.
	\item Ukuran naskah UNESCO (lebar 15,5 cm $\times$ tinggi 23 cm), Times New Roman 11 pt, spasi 1,5 (spasi 1 untuk abstrak, isi tabel, dan daftar pustaka), rata kiri-kanan (\textit{justify}). Margin atas 2 cm, kiri 2,5 cm, bawah 2 cm, kanan 2 cm; dicetak bolak-balik dengan margin cermin (margin 2,5 cm selalu di sisi jilid).
	\item Setiap anggota menyusun Laporan TA individu atas namanya sendiri dengan judul berbeda, membahas bagian proyek yang menjadi tanggung jawabnya (Pedoman TA 5.1.3). Bagian konteks proyek (profil dan masalah mitra, gambaran dan arsitektur umum, tinjauan pustaka bersama) boleh sama antaranggota; tujuan, rancangan, implementasi, pengujian, pembahasan, dan kesimpulan bagian individu wajib ditulis sendiri. Judul laporan individu ditetapkan saat seminar proposal dan dicantumkan dalam Surat Pembagian Kerja Tim.
	\item Proyek dikerjakan kelompok 3 sampai 4 mahasiswa bersama mitra berbadan hukum. Tiap anggota menyimpan \textit{logbook} individu dan mengisi \textit{peer assessment} sendiri secara rahasia; evaluasi kontribusi oleh pembimbing diisi terpisah (Pedoman TA butir 3.3.2 dan 3.3.4).
	\item Proyek dikerjakan kelompok 3 sampai 4 mahasiswa (lebih dari 4 dengan persetujuan Kaprodi) bersama mitra berbadan hukum, dengan TKT minimal 3 (Pedoman TA butir 3.3.2 dan 3.3.3).
	\item Seminar proposal dilaksanakan sebelum pelaksanaan dan bersamaan untuk seluruh anggota; sidang akhir dilaksanakan setelah luaran dievaluasi mitra.
	\item Dokumen mitra (surat kesediaan, evaluasi atau serah terima) ditandatangani \textit{Top Management} mitra.
	\item Halaman Pernyataan Penggunaan AI WAJIB diisi tiap anggota pada sidang akhir, termasuk bila tidak memakai AI (Pedoman TA 4.2.1).
	\item \textit{Code walkthrough} WAJIB pada sidang akhir bila luaran memuat perangkat lunak; setiap anggota menjelaskan kode bagiannya. Siapkan peta kode pada Lampiran 9.
	\item Batas pengerjaan 3 semester aktif sejak seminar proposal; bila tidak selesai, kelompok wajib mengajukan topik baru. Penanganan anggota bermasalah, anggota tersisa kurang dari 3, dan mitra yang tidak hadir mengikuti Pedoman TA butir 3.3.3 dan 3.3.5.
	\item Jalur kompetisi (juara 1 sampai 3 BELMAWA) hanya berlaku bila relevan; lampirkan sertifikat, dokumentasi, dan surat keterangan (Pedoman TA butir 3.3.6).
	\item \textit{Similarity} $\leq$25\% dan kutipan langsung $\leq$10\%, diuji per laporan; kemiripan dengan laporan anggota tim pada bagian konteks proyek boleh dikecualikan (Pedoman TA 5.7). Sitasi gaya IEEE.
\end{itemize}

\textbf{Cara penilaian}

\textbf{Seminar proposal.} Tabel 5.7 pedoman. Kualitas proposal dinilai dengan rubrik skema: kesiapan kemitraan (syarat), masalah mitra dan tujuan 30, rancangan dan rencana uji 30, kesiapan pelaksanaan 20, penguasaan rencana 20.

\textbf{Syarat masuk sidang.} Lulus/tidak tanpa skor: 3--4 anggota, mitra berbadan hukum, surat kesediaan, kehadiran \textit{Top Management} atau Form Selesai TA, TKT 3, \textit{logbook}, \textit{peer assessment}, batas 3 semester aktif.

\textbf{Sidang akhir.} Tabel 3.6: kualitas produk 50\%, kontribusi individu 30\%, presentasi individu 20\%. Kualitas produk dinilai dengan rubrik skema: penerimaan mitra 20 (ditolak mitra = tidak lulus), keterpenuhan tujuan 20, demonstrasi 10.
\clearpage
